\documentclass[12pt,a4paper]{article}

\usepackage[indonesian]{babel}
\usepackage[utf8]{inputenc}
\usepackage[T1]{fontenc}

\usepackage{geometry}
\geometry{
left=3cm,
right=3cm,
top=3cm,
bottom=3cm
}

\usepackage{newtxtext, newtxmath}
\usepackage{setspace}
\onehalfspacing

\usepackage{graphicx}
\usepackage{booktabs}
\usepackage{longtable}
\usepackage{array}
\usepackage{float}
\usepackage{hyperref}

\usepackage{tikz}
\usepackage{forest}

\usetikzlibrary{positioning}

\begin{document}

\begin{titlepage}
    \centering
    
    {\Large \textbf{IDENTIFIKASI PROYEK PERANGKAT LUNAK}} \\[0.4cm]
    {\large Pengembangan Aplikasi Absensi Mahasiswa Berbasis Mobile}

    {Disusun untuk memenuhi Tugas Mata Kuliah Manajemen Proyek Perangkat Lunak}\\
    {Dosen Pengampu: Yulis Nuryanti, S.Kom., M.Kom.}
    
    \vspace{2.5cm}
    
    \includegraphics[width=5cm]{images/unipi.png}
    
    \vspace{2.5cm} 
    
    {\large Disusun Oleh:} \\[0.4cm]
    \begin{tabular}{ll}
        \textbf{Nama} & : Wardiyansyah \\
        \textbf{NPM}  & : 2024807039 \\
    \end{tabular} \\[0.6cm]
    
    \vspace{\fill}
    {\large \textbf{PROGRAM STUDI SOFTWARE ENGINEERING}}\\
    {\large \textbf{FAKULTAS ILMU KOMPUTER}}\\
    {\large \textbf{UNIVERSITAS INSAN PEMBANGUNAN INDONESIA}} \\
    {\large \textbf{2026}}
    
\end{titlepage}
% ==================== END OF COVER ====================

\newpage

\tableofcontents

\newpage

%%%%%%%%%%%%%%%%%%%%%%%%%%%%%%%%%%%%%%%%%%%%%%%%%%

\section{Pendahuluan}

\subsection{Latar Belakang}

Perkembangan teknologi informasi mendorong berbagai institusi pendidikan untuk melakukan transformasi digital. Salah satu aktivitas yang masih banyak dilakukan secara manual adalah proses absensi mahasiswa.

Absensi manual memiliki beberapa kelemahan seperti membutuhkan waktu yang lama, rawan kesalahan pencatatan, sulit direkap, serta berpotensi terjadi manipulasi data. Oleh karena itu diperlukan suatu sistem absensi berbasis mobile yang mampu meningkatkan efektivitas dan efisiensi proses pencatatan kehadiran.

Menurut PMBOK, proyek merupakan suatu usaha sementara yang dilakukan untuk menghasilkan produk atau layanan yang unik. Pengembangan aplikasi absensi mahasiswa memenuhi karakteristik tersebut karena memiliki tujuan yang jelas, waktu pelaksanaan tertentu, anggaran, sumber daya manusia, serta menghasilkan sebuah perangkat lunak yang unik.

%%%%%%%%%%%%%%%%%%%%%%%%%%%%%%%%%%%%%%%%%%%%%%%%%%

\subsection{Rumusan Masalah}

\begin{enumerate}
\item Apa saja kegiatan dalam proyek pengembangan aplikasi absensi?
\item SDM apa saja yang dibutuhkan?
\item Keterampilan apa yang diperlukan?
\item Berapa estimasi waktu pengerjaan proyek?
\item Berapa estimasi biaya proyek?
\end{enumerate}

%%%%%%%%%%%%%%%%%%%%%%%%%%%%%%%%%%%%%%%%%%%%%%%%%%

\subsection{Tujuan}

\begin{enumerate}
\item Mengidentifikasi aktivitas proyek.
\item Mengidentifikasi kebutuhan SDM.
\item Menentukan keterampilan yang diperlukan.
\item Mengestimasi waktu proyek.
\item Mengestimasi biaya proyek.
\end{enumerate}

%%%%%%%%%%%%%%%%%%%%%%%%%%%%%%%%%%%%%%%%%%%%%%%%%%

\section{Landasan Teori}

\subsection{Pengertian Proyek}

Menurut PMBOK, proyek adalah usaha sementara yang dilakukan untuk menghasilkan produk, layanan, atau hasil yang unik.

Suatu proyek memiliki karakteristik:

\begin{itemize}
\item Memiliki awal dan akhir.
\item Bersifat sementara.
\item Memiliki ruang lingkup tertentu.
\item Memiliki biaya.
\item Memiliki jadwal.
\item Menghasilkan produk yang unik.
\end{itemize}

%%%%%%%%%%%%%%%%%%%%%%%%%%%%%%%%%%%%%%%%%%%%%%%%%%

\subsection{Manajemen Proyek}

Manajemen proyek merupakan proses penerapan pengetahuan, keterampilan, alat, dan teknik dalam melaksanakan aktivitas proyek agar tujuan proyek dapat dicapai sesuai ruang lingkup, waktu, biaya, dan kualitas yang telah ditentukan.

%%%%%%%%%%%%%%%%%%%%%%%%%%%%%%%%%%%%%%%%%%%%%%%%%%

\subsection{Triple Constraint}

Keberhasilan proyek dipengaruhi oleh tiga faktor utama yaitu:

\begin{enumerate}
\item Scope
\item Time
\item Cost
\end{enumerate}

Ketiga aspek tersebut harus dikelola secara seimbang.

%%%%%%%%%%%%%%%%%%%%%%%%%%%%%%%%%%%%%%%%%%%%%%%%%%

\section{Identifikasi Proyek}

\subsection{Nama Proyek}

Pengembangan Aplikasi Absensi Mahasiswa Berbasis Mobile.

%%%%%%%%%%%%%%%%%%%%%%%%%%%%%%%%%%%%%%%%%%%%%%%%%%

\subsection{Deskripsi Proyek}

Proyek ini bertujuan membangun aplikasi mobile yang digunakan oleh mahasiswa dan dosen dalam melakukan proses absensi secara digital.

Fitur utama aplikasi:

\begin{itemize}
\item Login
\item Absensi berbasis lokasi
\item Riwayat absensi
\item Rekapitulasi kehadiran
\item Dashboard dosen
\item Laporan kehadiran
\end{itemize}

%%%%%%%%%%%%%%%%%%%%%%%%%%%%%%%%%%%%%%%%%%%%%%%%%%

\subsection{Tujuan Proyek}

\begin{enumerate}
\item Mempermudah proses absensi.
\item Mengurangi penggunaan kertas.
\item Mempercepat rekapitulasi data.
\item Meningkatkan akurasi data.
\end{enumerate}

%%%%%%%%%%%%%%%%%%%%%%%%%%%%%%%%%%%%%%%%%%%%%%%%%%

\subsection{Ruang Lingkup}

Ruang lingkup proyek meliputi:

\begin{itemize}
\item Analisis kebutuhan
\item Desain sistem
\item Pembuatan database
\item Backend API
\item Aplikasi Android
\item Pengujian
\item Implementasi
\end{itemize}

\section{Studi Kelayakan Proyek}

\subsection{Kelayakan Teknis}

Dari sisi teknis, proyek ini layak dilaksanakan karena teknologi yang digunakan telah banyak diterapkan dalam pengembangan aplikasi mobile.

Teknologi yang digunakan:
\begin{itemize}
    \item Flutter
    \item Laravel
    \item MySQL
    \item Android Studio
    \item Visual Studio Code
    \item GitHub
\end{itemize}
Seluruh perangkat lunak tersebut bersifat open source sehingga tidak memerlukan biaya lisensi.

\subsection{Kelayakan Operasional}

Sistem dapat digunakan oleh seluruh dosen dan mahasiswa karena hampir seluruh civitas akademika telah memiliki smartphone Android.

Proses absensi menjadi lebih cepat dibandingkan absensi menggunakan kertas sehingga meningkatkan efisiensi pekerjaan bagian akademik.

\subsection{Kelayakan Ekonomi}

Penerapan sistem mampu mengurangi biaya operasional berupa:

\begin{itemize}
    \item kertas absensi
    \item Penyimpanan dokumen
    \item Waktu rekap absensi
    \item Kesalahan pencatatan data
\end{itemize}

Walaupun memerlukan biaya pengembangan pada awal proyek, dalam jangka panjang sistem mampu menghemat biaya operasional kampus.

\subsection{Kelayakan Hukum}

Pengembangan aplikasi mengikuti ketentuan mengenai perlindungan data pribadi serta hanya menyimpan informasi yang berkaitan dengan kegiatan akademik.

Hak akses dibedakan menjadi:
\begin{itemize}
    \item Administrator
    \item dosen
    \item Mahasiswa
\end{itemize}

\subsection{Kelayakan Jadwal}

Estimasi waktu pengembangan adalah 12 minggu, sehingga proyek dinilai layak dilaksanakan oleh tim yang terdiri dari enam sampai tujuh orang.

%%%%%%%%%%%%%%%%%%%%%%%%%%%%%%%%%%%%%%%%%%%%%%%%%%

\section{Perencanaan Proyek}

\subsection{Stakeholder}

Stakeholder yang terlibat pada proyek ini terdiri atas:

\begin{table}[H]
\centering
\begin{tabular}{ll}
\toprule
Stakeholder & Peran\\
\midrule
Project Owner & Menyetujui proyek dan menyediakan anggaran\\
Project Manager & Mengelola keseluruhan proyek\\
Dosen & Memberikan kebutuhan sistem\\
Mahasiswa & Pengguna Aplikasi\\
Tim Developer & Mengembangkan aplikasi\\
QA Tester & Menguji aplikasi\\
Bagian Akademik & Mengelola data absensi\\
\bottomrule
\end{tabular}
\end{table}

%%%%%%%%%%%%%%%%%%%%%%%%%%%%%%%%%%%%%%%%%%%%%%%%%%

\subsection{Sumber Daya Manusia}

\begin{table}[H]
\centering
\begin{tabular}{lc}
\toprule
Jabatan & Jumlah\\
\midrule
Project Manager & 1\\
System Analyst & 1\\
UI/UX Designer & 1\\
Backend Developer & 1\\
Mobile Developer & 1\\
QA Tester & 1\\
Trainer & 1\\
\bottomrule
\end{tabular}
\end{table}

%%%%%%%%%%%%%%%%%%%%%%%%%%%%%%%%%%%%%%%%%%%%%%%%%%

\subsection{Keterampilan yang Dibutuhkan}

\begin{table}[H]
    \centering
    \begin{tabular}{ll}
    \toprule
    Jabatan & Keterampilan\\
    \midrule
    Project Manager & Leadership, Scheduling, Risk Management\\
    System Analyst & Analisis kebutuhan sistem\\
    UI Designer & Figma, UI/UX Design\\
    Backend Developer & Laravel, REST API, PHP\\
    Mobile Developer & Flutter, Dart\\
    QA Tester & Black Box Testing, Bug Reporting\\
    Trainer & Public Speaking, Dokumentasi\\
    \bottomrule
    \end{tabular}
\end{table}

%%%%%%%%%%%%%%%%%%%%%%%%%%%%%%%%%%%%%%%%%%%%%%%%%%

\section{Work Breakdown Structure (WBS)}

Work Breakdown Structure (WBS) merupakan proses penguraian pekerjaan proyek menjadi beberapa bagian yang lebih kecil sehingga lebih mudah direncanakan, dikelola, dan dikendalikan. WBS pada proyek Pengembangan Aplikasi Absensi Mahasiswa Berbasis Mobile ditunjukkan pada Gambar.

\begin{figure}[H]
\centering
\resizebox{\textwidth}{!}{
    \begin{forest}
        for tree={
            draw,
            rounded corners,
            align=center,
            parent anchor=south,
child anchor=north,
edge={-latex},
l sep=12pt,
s sep=8pt
}
[Pengembangan\\Aplikasi Absensi
[Inisiasi
[Identifikasi Masalah]
[Studi Kelayakan]
[Penyusunan Proposal]
]
[Analisis
[Wawancara]
[Analisis Kebutuhan]
[Dokumen SRS]
]
[Perancangan
[Desain Database]
[Desain UI]
[Desain API]
]
[Pengembangan
[Backend]
[Mobile]
[Integrasi]
]
[Pengujian
[Unit Testing]
[Integration Testing]
[User Acceptance Test]
]
[Implementasi
[Deploy Server]
[Pelatihan User]
]
    [Penutupan
    [Dokumentasi]
    [Evaluasi]
    ]
    ]
    \end{forest}
}
\end{figure}

\newpage

\subsection{Rincian Work Breakdown Structure}

\begin{longtable}{cp{5cm}p{7cm}}
    \toprule
    \textbf{Kode} & \textbf{Pekerjaan} & \textbf{Deskripsi} \\
\midrule
1 & Inisiasi & Menentukan ruang lingkup proyek dan studi kelayakan. \\
2 & Analisis & Mengumpulkan kebutuhan pengguna dan menyusun Software Requirement Specification (SRS). \\
3 & Perancangan & Mendesain database, antarmuka pengguna, dan API sistem. \\
4 & Pengembangan & Mengembangkan backend, aplikasi mobile, dan integrasi sistem. \\
5 & Pengujian & Melakukan pengujian aplikasi sebelum implementasi. \\
6 & Implementasi & Deploy aplikasi ke server dan memberikan pelatihan kepada pengguna. \\
7 & Penutupan & Dokumentasi proyek dan evaluasi hasil pekerjaan. \\
\bottomrule
\end{longtable}

%%%%%%%%%%%%%%%%%%%%%%%%%%%%%%%%%%%%%%%%%%%%%%%%%%

\section{Penjadwalan Proyek}

\subsection{Jenis Pekerjaan}

Berdasarkan Work Breakdown Structure, aktivitas proyek dijabarkan menjadi beberapa pekerjaan utama yang akan digunakan dalam penyusunan jadwal proyek.

\begin{table}[H]
\centering
\begin{tabular}{cp{6cm}cc}
\toprule
\textbf{Kode} & \textbf{Aktivitas} & \textbf{Pendahulu} & \textbf{Durasi (Hari)} \\
\midrule
A & Inisiasi Proyek & - & 3 \\
B & Analisis Kebutuhan & A & 5 \\
C & Perancangan Sistem & B & 4 \\
D & Desain Database & C & 3 \\
E & Pengembangan Backend & D & 8 \\
F & Pengembangan Mobile & D & 8 \\
G & Integrasi Sistem & E,F & 3 \\
H & Pengujian Sistem & G & 5 \\
I & Implementasi & H & 2 \\
J & Penutupan & I & 1 \\
\bottomrule
\end{tabular}
\end{table}

\newpage    
\subsection{Diagram Network}

Diagram jaringan proyek ditunjukkan pada Gambar.

\begin{figure}[H]
\centering
\begin{tikzpicture}[>=latex,node distance=1.8cm]

\tikzstyle{activity}=[circle,draw,minimum size=9mm]

\node[activity](A){A};
\node[activity](B)[below=of A]{B};
\node[activity](C)[below=of B]{C};
\node[activity](D)[below=of C]{D};

\node[activity](E)[below left=1.5cm and 1.2cm of D]{E};
\node[activity](F)[below right=1.5cm and 1.2cm of D]{F};

\node[activity](G)[below=3cm of D]{G};
\node[activity](H)[below=of G]{H};
\node[activity](I)[below=of H]{I};
\node[activity](J)[below=of I]{J};

\draw[->](A)--(B);
\draw[->](B)--(C);
\draw[->](C)--(D);

\draw[->](D)--(E);
\draw[->](D)--(F);

\draw[->](E)--(G);
\draw[->](F)--(G);

\draw[->](G)--(H);
\draw[->](H)--(I);
\draw[->](I)--(J);

\end{tikzpicture}
\end{figure}

\subsection{Perhitungan Jalur Kritis (CPM)}

\begin{table}[H]
\centering
\begin{tabular}{ccccccc}
\toprule
Aktivitas & Durasi & ES & EF & LS & LF & Slack \\
\midrule
A & 3 & 0 & 3 & 0 & 3 & 0 \\
B & 5 & 3 & 8 & 3 & 8 & 0 \\
C & 4 & 8 & 12 & 8 & 12 & 0 \\
D & 3 & 12 & 15 & 12 & 15 & 0 \\
E & 8 & 15 & 23 & 15 & 23 & 0 \\
F & 8 & 15 & 23 & 15 & 23 & 0 \\
G & 3 & 23 & 26 & 23 & 26 & 0 \\
H & 5 & 26 & 31 & 26 & 31 & 0 \\
I & 2 & 31 & 33 & 31 & 33 & 0 \\
J & 1 & 33 & 34 & 33 & 34 & 0 \\
\bottomrule
\end{tabular}
\end{table}

\subsection{Jalur Kritis}

Berdasarkan hasil perhitungan CPM, jalur kritis proyek adalah:

\[
A \rightarrow B \rightarrow C \rightarrow D \rightarrow E \rightarrow G \rightarrow H \rightarrow I \rightarrow J
\]

Durasi penyelesaian proyek adalah \textbf{34 hari kerja}. Seluruh aktivitas pada jalur kritis memiliki nilai \textit{slack} sebesar nol, sehingga keterlambatan pada salah satu aktivitas tersebut akan mengakibatkan keterlambatan penyelesaian proyek secara keseluruhan.

\subsection{Gantt Chart}

Tabel berikut menunjukkan jadwal pelaksanaan proyek selama 34 hari kerja.

\begin{table}[H]
\centering
\begin{tabular}{clc}
\toprule
Kode & Aktivitas & Hari \\
\midrule
A & Inisiasi Proyek & 1--3 \\
B & Analisis Kebutuhan & 4--8 \\
C & Perancangan Sistem & 9--12 \\
D & Desain Database & 13--15 \\
E & Pengembangan Backend & 16--23 \\
F & Pengembangan Mobile & 16--23 \\
G & Integrasi Sistem & 24--26 \\
H & Pengujian Sistem & 27--31 \\
I & Implementasi & 32--33 \\
J & Penutupan & 34 \\
\bottomrule
\end{tabular}
\end{table}

%%%%%%%%%%%%%%%%%%%%%%%%%%%%%%%%%%%%%%%%%%%%%%%%%%

\section{Manajemen Risiko}

\subsection{Identifikasi Risiko}

Manajemen risiko dilakukan untuk mengidentifikasi berbagai kemungkinan yang dapat menghambat keberhasilan proyek. Dengan melakukan identifikasi risiko sejak awal, tim proyek dapat menyiapkan strategi mitigasi sehingga dampak yang ditimbulkan dapat diminimalkan.

Risiko yang diidentifikasi pada proyek Pengembangan Aplikasi Absensi Mahasiswa Berbasis Mobile ditunjukkan pada Tabel.

\begin{table}[H]
\centering
\begin{tabular}{p{0.8cm} p{4cm} p{2.3cm} p{2cm} p{5cm}}
\toprule
\textbf{No} & \textbf{Risiko} & \textbf{Kemungkinan} & \textbf{Dampak} & \textbf{Mitigasi} \\
\midrule
1 & Perubahan kebutuhan pengguna & Sedang & Tinggi & Melakukan validasi kebutuhan dan persetujuan dokumen SRS sebelum pengembangan dimulai. \\

2 & Keterlambatan pengembangan & Sedang & Tinggi & Monitoring progres mingguan dan evaluasi pencapaian setiap sprint. \\

3 & Anggota tim sakit atau mengundurkan diri & Rendah & Sedang & Menyiapkan dokumentasi pekerjaan dan melakukan pembagian tugas kepada anggota lain. \\

4 & Server mengalami gangguan & Sedang & Tinggi & Menyediakan server cadangan serta melakukan backup data secara berkala. \\

5 & Bug pada aplikasi saat implementasi & Tinggi & Sedang & Melaksanakan unit testing, integration testing, dan user acceptance testing sebelum implementasi. \\

6 & Kehilangan data proyek & Rendah & Tinggi & Melakukan backup repository GitHub dan database setiap hari. \\
\bottomrule
\end{tabular}
\end{table}

\subsection{Analisis Risiko}

Berdasarkan identifikasi risiko yang telah dilakukan, risiko dengan prioritas tertinggi adalah perubahan kebutuhan sistem, keterlambatan pengembangan, dan gangguan server. Risiko-risiko tersebut memiliki dampak yang besar terhadap keberhasilan proyek sehingga perlu dilakukan pengawasan secara berkala oleh Project Manager.

\subsection{Strategi Penanganan Risiko}

Strategi yang digunakan dalam proyek ini meliputi:

\begin{enumerate}
\item Melakukan rapat koordinasi proyek setiap minggu.
\item Menyusun dokumen kebutuhan sistem secara lengkap sebelum tahap pengembangan.
\item Menggunakan sistem version control (Git) untuk menghindari kehilangan kode program.
\item Melakukan backup database dan source code secara berkala.
\item Melaksanakan pengujian aplikasi sebelum implementasi kepada pengguna.
\item Menyediakan waktu cadangan (buffer) pada jadwal proyek untuk mengantisipasi keterlambatan pekerjaan.
\end{enumerate}

%%%%%%%%%%%%%%%%%%%%%%%%%%%%%%%%%%%%%%%%%%%%%%%%%%

\section{Manajemen Biaya}

\subsection{Estimasi Biaya Proyek}

Estimasi biaya dilakukan untuk mengetahui besarnya anggaran yang diperlukan selama pelaksanaan proyek. Biaya proyek terdiri atas biaya sumber daya manusia, biaya operasional, dan cadangan risiko.

\subsection{Biaya Sumber Daya Manusia}

\begin{table}[H]
\centering
\begin{tabular}{p{0.8cm}p{4cm}r}
\toprule
\textbf{No} & \textbf{Jabatan} & \textbf{Biaya (Rp)}\\
\midrule
1 & Project Manager & 15.000.000\\
2 & System Analyst & 8.000.000\\
3 & UI/UX Designer & 3.500.000\\
4 & Backend Developer & 13.500.000\\
5 & Mobile Developer & 13.500.000\\
6 & Quality Assurance & 3.000.000\\
7 & Trainer & 2.000.000\\
\midrule
\multicolumn{2}{r}{\textbf{Subtotal}} & \textbf{58.500.000}\\
\bottomrule
\end{tabular}
\end{table}

\subsection{Biaya Operasional}

\begin{table}[H]
\centering
\begin{tabular}{p{0.8cm}p{6cm}r}
\toprule
\textbf{No} & \textbf{Kebutuhan} & \textbf{Biaya (Rp)}\\
\midrule
1 & Hosting Server & 2.000.000\\
2 & Domain & 500.000\\
3 & Internet dan Komunikasi & 1.500.000\\
4 & Konsumsi Meeting & 1.500.000\\
5 & Dokumentasi & 1.000.000\\
6 & Peralatan Pendukung & 5.000.000\\
\midrule
\multicolumn{2}{r}{\textbf{Subtotal}} & \textbf{11.500.000}\\
\bottomrule
\end{tabular}
\end{table}

\subsection{Cadangan Risiko}

Cadangan risiko (\textit{contingency reserve}) dialokasikan sebesar 10\% dari total biaya proyek sebagai antisipasi terhadap perubahan kebutuhan, kenaikan harga perangkat, maupun keterlambatan pekerjaan.

\[
\text{Cadangan Risiko}
=
10\%
\times
(Rp58.500.000 + Rp11.500.000)
=
Rp7.000.000
\]

\subsection{Total Anggaran Proyek}

\begin{table}[H]
\centering
\begin{tabular}{lr}
\toprule
\textbf{Komponen} & \textbf{Biaya (Rp)}\\
\midrule
Biaya SDM & 58.500.000\\
Biaya Operasional & 11.500.000\\
Cadangan Risiko & 7.000.000\\
\midrule
\textbf{Total Anggaran Proyek} & \textbf{77.000.000}\\
\bottomrule
\end{tabular}
\end{table}

\subsection{Sumber Dana}

Proyek ini direncanakan menggunakan kombinasi beberapa sumber pendanaan agar pelaksanaan proyek dapat berjalan dengan baik. Rincian sumber dana disajikan pada Tabel.

\begin{table}[H]
\centering
\begin{tabular}{lr}
\toprule
\textbf{Sumber Dana} & \textbf{Persentase}\\
\midrule
Dana Universitas & 70\%\\
Hibah Pengembangan Sistem & 20\%\\
Kerja Sama/Sponsor & 10\%\\
\bottomrule
\end{tabular}
\end{table}

\subsection{Pengendalian Biaya}

Pengendalian biaya dilakukan selama pelaksanaan proyek agar pengeluaran tetap sesuai dengan anggaran yang telah direncanakan. Mekanisme pengendalian biaya meliputi:

\begin{enumerate}
\item Menyusun anggaran proyek sebelum pelaksanaan dimulai.
\item Mencatat seluruh pengeluaran proyek secara berkala.
\item Membandingkan biaya aktual dengan anggaran yang telah disusun.
\item Melakukan evaluasi apabila terjadi penyimpangan biaya lebih dari 10\%.
\item Menggunakan cadangan risiko hanya setelah memperoleh persetujuan dari Project Manager dan Project Owner.
\item Menyusun laporan biaya proyek pada setiap akhir bulan sebagai bahan evaluasi.
\end{enumerate}
%%%%%%%%%%%%%%%%%%%%%%%%%%%%%%%%%%%%%%%%%%%%%%%%%%

\section{Kesimpulan}

Berdasarkan hasil identifikasi dan perencanaan proyek, pengembangan aplikasi absensi mahasiswa berbasis mobile merupakan proyek perangkat lunak yang layak untuk dilaksanakan. Proyek ini memiliki tujuan, ruang lingkup, kebutuhan sumber daya manusia, jadwal pelaksanaan, serta estimasi biaya yang telah direncanakan secara sistematis. Melalui penyusunan Work Breakdown Structure (WBS), penjadwalan proyek menggunakan Critical Path Method (CPM), serta identifikasi jalur kritis, seluruh aktivitas proyek dapat direncanakan dan dikendalikan dengan lebih efektif.

Selain itu, penerapan manajemen risiko dan manajemen biaya memberikan pedoman dalam mengantisipasi berbagai kendala yang mungkin terjadi selama pelaksanaan proyek, sehingga penggunaan sumber daya dapat lebih efisien dan anggaran tetap terkendali. Berdasarkan hasil estimasi, proyek direncanakan memiliki durasi pelaksanaan selama 34 hari kerja dengan total anggaran sebesar Rp77.000.000. Dengan perencanaan yang baik, proyek ini diharapkan mampu menghasilkan aplikasi absensi yang efektif, meningkatkan efisiensi pengelolaan kehadiran mahasiswa, serta mendukung proses administrasi akademik secara lebih akurat dan terintegrasi.
\end{document}
